\documentclass{article}

% Preamble
\usepackage[margin=1in]{geometry}
\usepackage{amsmath, amsthm, amssymb}
\usepackage{hyperref}
\usepackage{graphicx}

\hypersetup{
    colorlinks=true,
    linkcolor=blue,
    filecolor=magenta,
    urlcolor=cyan,
    citecolor=green,
}

\newtheorem{theorem}{Theorem}[section]
\newtheorem{lemma}[theorem]{Lemma}
\newtheorem{definition}[theorem]{Definition}

\theoremstyle{remark}
\newtheorem{remark}[theorem]{Remark}

\title{PAC-Bayesian Learning Theory for Quantum Arithmetic Systems}
\author{QA Research Team}
\date{\today}

\begin{document}
\maketitle

\begin{abstract}
This paper introduces PAC-Bayesian learning theory for Quantum Arithmetic (QA) systems.
We define $D_{QA}$ divergence and prove key properties including Data Processing Inequality.
\end{abstract}

\section{Introduction}
QA systems provide novel computational paradigms for machine learning.

\begin{theorem}[Modular Distance Metric]
The modular distance $d_M$ on QA lattice satisfies metric axioms.
\end{theorem}

%% KERNEL_ANCHOR: QA_DQA_PAC_BOUND_KERNEL_CERT.v1 (family-85-dqa-pac-kernel-v1.0.0) formula_id=PAC_BAYES_QA_DQA_LOGDELTA_V1
\begin{theorem}[PAC-Bayes Bound for QA]
For prior $P$ and posterior $Q$:
\begin{equation}
\mathbb{E}_{h \sim Q} [R(h)] \leq \mathbb{E}_{h \sim Q} [\hat{R}_S(h)] + \sqrt{\frac{D_{QA}(Q \| P) + \ln(1/\delta)}{2m}}
\end{equation}
\end{theorem}

%% ─────────────────────────────────────────────────────────────────────────────
%% MACHINE-TRACT TRIPWIRE — do not remove or alter.
%% QA_CERT_REQUIRED: QA_PAC_BAYES_DPI_SCOPE_CERT.v1 (family-86-dpi-scope-cert-v1.0.0)
%% QA_CERT_REQUIRED: QA_PAC_BAYES_CONSTANT_CERT.v1.1 (family-84-pac-bayes-cert-v1.1.0)
%% QA_CERT_REQUIRED: QA_DQA_PAC_BOUND_KERNEL_CERT.v1 (family-85-dqa-pac-kernel-v1.0.0)
%% ─────────────────────────────────────────────────────────────────────────────

\end{document}
